% MedVision manuscript -- Discussion section

\section{Discussion}\label{sec:discussion}

\subsection{Principal findings}\label{sec:discussion-principal}
The central result was a mismatch between high predictive performance and
incremental image evidence. Gray-Square retained Macro-F1 0.9381 versus the
baseline eight-seed mean of $0.9397\pm0.0061$; baseline ICM was
$0.0016\pm0.0065$, with a non-significant signed test and TOST equivalence to
the operational interval $[-0.01,0.01]$. Thus, replacing the radiograph changed
little under this intervention. It does not imply that radiographs are
universally unused.

The tested interventions did not provide a consistent remedy. LMH had the
highest Macro-F1 but lower AUROC and substantially worse ECE and Brier score;
Noise-Consistency also increased Brier score without increasing ICM.
Adaptive-CADQ remained near baseline, and the small GACR-B increase was
exploratory because it was added after the initial plan. Extracted image gates
(0.42624, 0.49963, and 0.57336 by class) stayed near initialization rather than
collapsing. This is a stagnation/null observation, not proof that the image
branch was disconnected.

We refer to this pattern as modality inertia, or gate stagnation: the gate
neither collapses nor learns because the report pathway already supplies a
strong label-correlated signal, leaving insufficient gradient incentive for
image reliance.

\textbf{Label--report circularity and endpoint validity.}
MIMIC-CXR radiographs are paired with reports, CheXpert labels are generated
from report language, and the model receives the report’s \texttt{findings}
field \citep{johnson2019mimiccxr,irvin2019chexpert}. The task therefore does not
isolate image-based disease recognition. Text-only Macro-F1/AUROC were 0.8555/
0.9820, while Gray-Square reached 0.9381/0.9960. A constant image consequently
preserved most of the multimodal score, consistent with report-linked shortcut
potential. The control does not identify the phrases used or show that labels
are clinically wrong; it identifies a construct-validity risk in which the
target and input share an information source.

This linkage explains why text-debiasing can appear incoherent. Suppressing a
report shortcut may remove information rewarded by a report-derived label
without creating an independent visual signal, yielding lower discrimination or
calibration with no ICM gain. The null result should therefore not be attributed
solely to a failed fusion module.

The GACR-B result is especially informative. Even the grounding-aware
consistency intervention proposed by Sadanandan and Behzadan
\citep{sadanandan2026consistentdangerous} could not overcome the circularity:
with $\lambda=0.3$, its consistency penalty did not generate enough gradient to
counter the dominant report shortcut. This is a limitation of the label--input
relationship, not a claim that the proposed intervention is ineffective under
independent labels.

\subsection{Architectural, statistical, and regulatory interpretation}
The attention block contains one projected image token and one projected text
token. Its single-key attention cannot express word--region alignment, so its
weights and Grad-CAM examples are not causal grounding evidence. High Macro-F1,
near-zero ICM, and a non-collapsed gate can coexist when the image is redundant
with text or receives little optimisation incentive. The archived negative-KL
Noise-Consistency objective additionally rewards prediction divergence when
minimised; a conventional consistency claim requires a positive-KL rerun.

The Gray-Square control was a single retained run reused for eight ICM values,
and image-only checkpoint restoration was inconsistent. These asymmetries make
the controls useful diagnostics but understate control uncertainty and prevent
modality ranking. Confirmatory work should train matched controls at every seed
and propagate their variation. Seed-level pairing is appropriate because the
same 455 studies recur; pooled subject results are descriptive only.

\textbf{Implications for evaluation and future study design.}
Modality contribution should be audited directly rather than inferred from a
high multimodal score. Future studies should combine report-preserving image
ablation with independent or image-derived labels, masked or temporally
separated report fields, counterfactual text tests, deterministic source
provenance, and external cohorts. ICM should be validated under these settings,
not treated as a standalone grounding certificate. Reporting controls, label
provenance, checkpoint rules, calibration, and seed-level uncertainty together
is consistent with CLAIM and TRIPOD+AI transparency principles
\citep{mongan2024claim,collins2024tripodai}.

For regulatory-facing evaluation, ICM and gate-stagnation checks should be
treated as evidence-traceability tests rather than optional visualisations. A
high aggregate score does not establish that the declared image input affected
the decision, particularly when labels and reports share provenance. FDA good
machine-learning-practice principles and its real-world-evidence guidance
emphasise traceable data provenance, representative evaluation, uncertainty, and
evidence that supports the intended use \citep{fda2025gmlp,fda2025rwe}. In that
spirit, a submission for a multimodal medical device should report the
report-preserving ablation, matched seed-level uncertainty, calibration, gate
behaviour, and independent-label validation before making a visual-grounding
claim.

The findings are not a clinical recommendation to ignore radiographs. They are
limited to one report-linked cohort, task, preprocessing path, and checkpoint
archive, with known normalisation, loss-sign, small-minority, post-hoc, and
single-control limitations. The defensible conclusion is that these
interventions did not establish robust incremental image evidence in this
benchmark; corrected implementations, independent labels, matched controls,
and external validation are needed before broader claims.
